\documentclass[aps,pra,reprint,amsmath,amssymb,superscriptaddress,longbibliography]{revtex4-2}

\usepackage{booktabs}
\usepackage{bm}
\usepackage{hyperref}
\usepackage{microtype}
\usepackage{xcolor}

\newcommand{\Ffour}{F_{4}}
\newcommand{\ii}{\mathrm{i}}
\newcommand{\ee}{\mathrm{e}}
\newcommand{\per}{\operatorname{per}}
\newcommand{\hist}{\operatorname{hist}}
\newcommand{\evec}{\bm e}

\begin{document}

\title{Directional response jets of suppressed multiphoton events\\
in a four-mode Fourier interferometer}

\author{Hainan Zhao}
\email{hainzhao@gmail.com}
\affiliation{Independent Researcher}

\date{\today}

\begin{abstract}
Suppressed output events are established signatures of multiphoton
interference, but a zero measured at one operating point does not reveal how
the event responds to coherent implementation errors.  We resolve this
ambiguity for selected events of the four-mode Fourier interferometer by
expanding their probabilities along a complete set of calibrated
$SU(4)$ output directions.  Depending on the event and direction, the
probability remains identically zero or begins at quadratic or quartic order.
For the four-photon transition
$(0,1,2,1)\rightarrow(0,1,2,1)$, one beam-splitter quadrature obeys
$P_{Y_{13}}(\epsilon)=0$ for every angle, whereas the orthogonal quadrature
obeys the exact law
$P_{X_{13}}(\epsilon)=\sin^2(2\epsilon)/16$.
The exact null extends to
$(0,a,2a,a)\rightarrow(0,a,2a,a)$ for every positive odd $a$.
As a structural tool, we derive an $F_4$-specific, phase-resolved sectorwise
reciprocity from classical binary Krawtchouk duality.  It maps this apparently
nonperiodic family to cyclic suppression and preserves reciprocal amplitudes
throughout the standard dephased order-four complex-Hadamard family.
We give all twelve off-diagonal $SU(4)$ response directions, an
eleven-photon example with quartic onset, and a finite-shot nuisance-floor
model.  These event-level response jets refine disorder-averaged robustness
estimates and provide falsifiable quadrature-sensitive checks for programmable
photonic circuits.
\end{abstract}

\maketitle

\section{Introduction}

Exact zeros in multiparticle transition probabilities generalize the
Hong--Ou--Mandel effect to multiport interferometers.  Fourier suppression
laws \cite{Tichy2010,Crespi2016,Dittel2018PRL,Dittel2018PRA} and their
extensions \cite{Bezerra2023} are useful both as interference phenomena and as
device diagnostics.  Experiments have used forbidden-event populations to
quantify departures from ideal behavior
\cite{Crespi2016,Wang2023}, while recent protocols use Fourier-suppressed
outputs to benchmark genuine multiphoton indistinguishability
\cite{Sanz2026}.  Partial distinguishability and generic weak disorder
therefore have a mature theory
\cite{Shchesnovich2015,Dittel2018PRL}.

Here we ask a different question: \emph{what is the coherent directional
response of a specified dark event?}  A random scalar disorder strength
averages over physically distinct tangent directions.  A programmable
interferometer instead allows one to append a calibrated two-mode rotation
and scan its signed angle.  The resulting Taylor series---the event's response
jet---can contain directional information that is erased by averaging.

Our principal results are:
\begin{enumerate}
\item exact neighboring-permanent formulas give the response along all twelve
off-diagonal generators of $SU(4)$, while the three diagonal directions
preserve every dark event;
\item selected Fourier-dark events exhibit exact, quadratic, and quartic
probability onsets, rather than one universal robustness exponent;
\item the four-photon event
$\bm r=\bm s=(0,1,2,1)$ has the exact quadrature contrast
\begin{equation}
 P_{Y_{13}}(\epsilon)=0,\qquad
 P_{X_{13}}(\epsilon)=\frac{\sin^2(2\epsilon)}{16};
\label{eq:headline}
\end{equation}
\item an $F_4$-specific sectorwise refinement of classical
Krawtchouk/symmetric-power reciprocity
\cite{Feinsilver2005,Chami2014,ChenLouck1998} organizes zeros that evade a
direct cyclic test.
\end{enumerate}

We do not claim that the displayed response of one event is a universal
invariant of every zero generated by the same broad suppression mechanism.
The experimentally falsifiable objects here are \emph{event--direction
pairs}.  This restriction is important: different members of a suppression
family can have different stabilizers and different response jets.

\section{Fourier amplitudes and phase bins}
\label{sec:setup}

Let $\bm r=(r_0,\ldots,r_3)$ and $\bm s=(s_0,\ldots,s_3)$ be input and
output occupations with $\sum_jr_j=\sum_js_j=n$.  For a single-particle
unitary $U$, the Fock transition amplitude is
\begin{equation}
\mathcal A_{\bm r,\bm s}(U)=
\frac{\per U[\bm s,\bm r]}
{\sqrt{\prod_jr_j!\prod_ks_k!}},
\label{eq:fock-amplitude}
\end{equation}
where rows and columns are repeated according to $\bm s$ and $\bm r$.
We use the canonical dephased and ordered Fourier matrix
\begin{equation}
(\Ffour)_{jk}=\frac{\ii^{jk}}{2},\qquad j,k\in\{0,1,2,3\}.
\end{equation}

Before the common factor $2^{-n}$ is inserted, each labelled permanent path
has one of four phases.  We write
\begin{equation}
\bm h(\bm r,\bm s)=(h_0,h_1,h_2,h_3),
\quad
\per (2\Ffour)[\bm s,\bm r]
=\sum_{\ell=0}^{3}h_\ell\ii^\ell .
\label{eq:hist}
\end{equation}
The histogram is a combinatorial object in this fixed gauge; it is not itself
invariant under arbitrary row and column rephasings.  Equality of transition
probabilities and simultaneous zero/nonzero status are gauge invariant.

For reference, the elementary cyclic rule follows by shifting a periodic
occupation.  If the input is invariant under a shift $\tau$, then a nonzero
transition requires
\begin{equation}
\tau\sum_{j=0}^{3}j s_j=0\pmod 4 ,
\label{eq:cyclic}
\end{equation}
with the reciprocal output condition following from the symmetry of
$\Ffour$ \cite{Tichy2010,Dittel2018PRA}.

\section{Phase-resolved sectorwise reciprocity}
\label{sec:reciprocity}

Classical symmetric-power identities transpose all input and output indices
for a symmetric matrix \cite{ChenLouck1998,Chami2014}.  The binary
factorization of $\Ffour$ permits a partial exchange: the odd-mode split can
be transposed while both even-mode splits are held fixed.

\paragraph*{Sectorwise reciprocity theorem.}
Let $d,N\geq0$, $0\leq\alpha,s\leq d$, and $0\leq k,p\leq N$.  Define
\begin{align}
\mathcal H_{d,N}(\alpha,k;s,p)
=\bm h\big(&(\alpha,k,d-\alpha,N-k),\nonumber\\
           &(s,p,d-s,N-p)\big).
\end{align}
Then
\begin{align}
\mathcal H(\alpha,k;s,p)
&=\mathcal H(\alpha,p;s,k)\nonumber\\
&=\mathcal H(s,k;\alpha,p)
=\mathcal H(s,p;\alpha,k).
\label{eq:fourfold}
\end{align}
Thus the complete four-bin phase count, not only the Fourier amplitude, is
preserved.  The visible hypothesis is that input and output have the same
total occupation $d$ in the even-indexed modes (and hence the same total
$N$ in the odd-indexed modes).

We give a compact proof.  For $q^4=1$, let
\begin{equation}
\Phi_q(\bm r,\bm s)=
\left(\prod_jr_j!\right)
[\bm x^{\bm r}]
\prod_{u=0}^{3}\left(\sum_{v=0}^{3}q^{uv}x_v\right)^{s_u}.
\label{eq:phi}
\end{equation}
The four values $\Phi_1,\Phi_{\ii},\Phi_{-1},\Phi_{-\ii}$ are the
four-point discrete Fourier transform of $\bm h$.
At $q=\ii$, introduce
\begin{equation}
E=x_0+x_2,\quad T=x_0-x_2,\quad
S=x_1+x_3,\quad D=x_1-x_3.
\end{equation}
The four forms become $E+S,T+\ii D,E-S,T-\ii D$.
After expanding the even sector, every remaining summand contains
\begin{equation}
K_m(p;N)K_k(m;N),
\end{equation}
where
\begin{equation}
K_m(p;N)=\sum_t(-1)^t
\binom pt\binom{N-p}{m-t}
\end{equation}
is a binary Krawtchouk coefficient.  Self-duality gives
\begin{equation}
\binom NpK_m(p;N)=\binom NmK_p(m;N).
\label{eq:kraw-duality}
\end{equation}
Applied twice, Eq.~\eqref{eq:kraw-duality} exchanges $p$ and $k$
term by term.  The binomial ratio is canceled exactly by the
$k!(N-k)!$ and $p!(N-p)!$ factors in Eq.~\eqref{eq:phi}.
At $q=-1$, the polynomial is $(E+S)^d(E-S)^N$; the requested total
degree $d$ in $(x_0,x_2)$ forces the $E^dS^N$ contribution, and the same
factorial cancellation applies.  The $q=1$ value counts all labelled paths,
and $q=-\ii$ follows by conjugation.  Fourier inversion proves
Eq.~\eqref{eq:fourfold}.  Details are in the Supplemental Material.

The operation is generically not a mode relabelling.  For example, with
$d=5$, $N=11$, $(\alpha,s,p,k)=(1,0,3,2)$, the input occupation multisets
on the two sides are $\{1,2,4,9\}$ and $\{1,3,4,8\}$, while the output
multisets are $\{0,3,5,8\}$ and $\{0,2,5,9\}$.  They remain unequal after
global input--output exchange.

\subsection{Polynomial reciprocity for the order-four Hadamard family}

Consider
\begin{equation}
H(z)=
\begin{pmatrix}
1&1&1&1\\
1&z&-1&-z\\
1&-1&1&-1\\
1&-z&-1&z
\end{pmatrix}.
\label{eq:hz}
\end{equation}
Replacing $\ii$ by $z$ in the preceding coefficient proof adds a common
factor $z^{N-m}$.  The reciprocal permanents are therefore equal as
polynomials in $z$.  When $|z|=1$, $H(z)/2$ is the standard dephased
one-parameter order-four complex-Hadamard family
\cite{TadejZyczkowski2006}, and the normalized Fock amplitudes are equal.
This is amplitude reciprocity, not persistence of every Fourier zero:
already for $\bm r=\bm s=(0,1,2,1)$ the relevant unnormalized coefficient
is proportional to $2(1+z^2)$ and is generically nonzero.

\subsection{Closure of a nominally noncyclic family}

For
\begin{equation}
\bm r=\bm s=(0,a,2a,a),
\label{eq:odd-family}
\end{equation}
rotations and Eq.~\eqref{eq:fourfold} transport the zero condition to
\begin{equation}
(0,0,2a,2a)\longrightarrow(a,a,a,a).
\end{equation}
The latter is suppressed by Eq.~\eqref{eq:cyclic} precisely when $a$ is
odd.  Thus Eq.~\eqref{eq:odd-family} is not evidence for a separate
selection mechanism; it lies in the reciprocity closure of a cyclic law.
An exact enumeration illustrates the compression produced by this closure:
\begin{center}
\begin{tabular}{crrr}
\toprule
$n$ & residual reps. & closure components & reach cyclic\\
\midrule
4&3&3&1\\
5&8&3&0\\
6&10&6&0\\
7&0&0&0\\
8&33&23&1\\
9&72&40&0\\
\bottomrule
\end{tabular}
\end{center}
These finite values are an application of the theorem, not an asymptotic
classification claim.

\section{Directional response jets}
\label{sec:jets}

Append a calibrated two-mode rotation after $\Ffour$:
\begin{align}
U^X_{pq}(\epsilon)&=\ee^{\ii\epsilon X_{pq}}\Ffour,
&X_{pq}&=|p\rangle\langle q|+|q\rangle\langle p|,\nonumber\\
U^Y_{pq}(\epsilon)&=\ee^{\ii\epsilon Y_{pq}}\Ffour,
&Y_{pq}&=-\ii|p\rangle\langle q|+\ii|q\rangle\langle p|.
\label{eq:mixers}
\end{align}
The six unordered pairs and two quadratures span the twelve off-diagonal
directions of $\mathfrak{su}(4)$.  The remaining three traceless diagonal
generators append output phases; they multiply a fixed Fock amplitude by a
phase and therefore preserve every exact zero.

Let $Z_{\bm r,\bm t}=\per(2\Ffour)[\bm t,\bm r]$ and
\begin{equation}
D_{\bm r,\bm s}=2^n
\sqrt{\prod_jr_j!\prod_ks_k!}.
\end{equation}
For a dark target, Fock ladder factors give the exact tangents
\begin{align}
\mathcal A'_X(0)
&=\frac{\ii}{D_{\bm r,\bm s}}
\left(s_qZ_{\bm r,\bm s+\evec_p-\evec_q}
+s_pZ_{\bm r,\bm s-\evec_p+\evec_q}\right),
\label{eq:x-tangent}\\
\mathcal A'_Y(0)
&=\frac{1}{D_{\bm r,\bm s}}
\left(s_pZ_{\bm r,\bm s-\evec_p+\evec_q}
-s_qZ_{\bm r,\bm s+\evec_p-\evec_q}\right).
\label{eq:y-tangent}
\end{align}
Negative occupations are omitted.  If the tangent is nonzero,
$P(\epsilon)=|\mathcal A'(0)|^2\epsilon^2+O(\epsilon^3)$.
If it vanishes, Krylov moments of the two-mode generator decide whether a
higher-order onset occurs or the amplitude is identically zero.

\subsection{Four-photon benchmark predictions}

Table~\ref{tab:four-photon} gives the complete off-diagonal response for two
specified four-photon dark events.  ``Exact'' means zero for every rotation
angle on that axis.
\begin{table}[b]
\caption{Leading ideal probabilities for two four-photon events.}
\label{tab:four-photon}
\begin{ruledtabular}
\begin{tabular}{c c c c c}
event & pair & $X_{pq}$ & $Y_{pq}$\\
\hline
cyclic &01&$3\epsilon^2/8$&$3\epsilon^2/8$\\
$\bm r=(1,1,1,1)$&02&exact&exact\\
$\bm s=(3,1,0,0)$&03&$3\epsilon^2/8$&$3\epsilon^2/8$\\
&12&exact&exact\\
&13&exact&exact\\
&23&exact&exact\\[2pt]
\hline
reciprocity-closed&01&$\epsilon^2/64$&$\epsilon^2/64$\\
$\bm r=\bm s=$&02&$\epsilon^2/4$&$\epsilon^2/4$\\
$(0,1,2,1)$&03&$\epsilon^2/64$&$\epsilon^2/64$\\
&12&$\epsilon^2/64$&$25\epsilon^2/64$\\
&13&$\epsilon^2/4$&exact\\
&23&$\epsilon^2/64$&$25\epsilon^2/64$\\
\end{tabular}
\end{ruledtabular}
\end{table}

For the cyclic example, arbitrary $U(2)$ rotations on pairs
$02,12,13,23$ preserve the zero.  Each rotation preserves the total
occupation in its pair, and none of the reachable modular sums satisfies
Eq.~\eqref{eq:cyclic}.  For the second example, direct coefficient
evaluation strengthens the small-angle $13$ entries to Eq.~\eqref{eq:headline}.
The two quadratures therefore give a clean same-pair contrast.

The $Y_{13}$ null is not a four-photon accident.  For every positive odd
$a$, the event in Eq.~\eqref{eq:odd-family} remains dark under the full
$Y_{13}$ rotation.  With $x=x_1$, $y=x_2$, $z=x_3$, $u=x+z$, and
$v=x-z$, its coefficient is proportional to
\begin{equation}
[x^ay^{2a}z^a]
\left(A(y^2+v^2)+Byv\right)^a(y-u)^{2a}.
\label{eq:protected-coefficient}
\end{equation}
Reflection $x\leftrightarrow z$ removes odd powers of $v$.  After choosing
$2h$ factors of $Byv$, the remaining sum pairs an index $k$ with $a-k$.
Their weights agree, whereas
\begin{equation}
\begin{aligned}
&[x^az^a](x-z)^{2(a-k)}(x+z)^{2k}\\
&\qquad=(-1)^a[x^az^a](x-z)^{2k}(x+z)^{2(a-k)}.
\end{aligned}
\end{equation}
Every pair cancels for odd $a$, for every $\epsilon$.

\subsection{A quartic direction}

Higher-order flatness need not imply exact protection.  For
\begin{equation}
\bm r=(0,1,3,7),\qquad \bm s=(1,3,3,4),
\end{equation}
an integral zero on a previously scanned affine family, the $Y_{12}$
tangent vanishes but the second derivative does not:
\begin{equation}
P^{Y_{12}}_{\bm r,\bm s}(\epsilon)
=\frac{315}{8192}\epsilon^4+O(\epsilon^5).
\label{eq:quartic}
\end{equation}
The other eleven off-diagonal directions and the exact Gaussian-integer
certificate are given in the Supplemental Material.  Equation
\eqref{eq:quartic} is a local quartic onset, not an all-angle null and not a
claim that the event is globally isolated among all occupations.

\section{Applications and finite-shot predictions}
\label{sec:applications}

\paragraph*{Quadrature-sensitive calibration.}
Equation~\eqref{eq:headline} predicts the entire ideal angle dependence for
two orthogonal quadratures on the same pair of modes.  Counts on the
$Y_{13}$ axis reveal departure from the joint ideal model, while the
$X_{13}$ curve supplies an in situ scale for the programmed angle.  This
does not replace full unitary tomography \cite{RahimiKeshari2013}; it is an
event-level, multiphoton consistency check.

\paragraph*{Cross-configuration consistency.}
Equation~\eqref{eq:fourfold} produces occupation pairs with equal ideal
probabilities for every $|z|=1$ in Eq.~\eqref{eq:hz}.  Sweeping $z$ and
comparing reciprocal configurations overdetermines the ideal prediction.
A discrepancy diagnoses failure of the joint state--interferometer--detector
model, but by itself does not identify which component failed.

\paragraph*{Protected analysis subgroups.}
The exact rows of Table~\ref{tab:four-photon} identify output-mode subgroups
within which a chosen target remains dark.  They can be used to modify the
analysis basis without opening that event.  This is an event-specific design
constraint, not a fault-tolerant subspace.

\paragraph*{Separating a floor from coherent response.}
To quantify a near-term experiment, we use the transparent nuisance model
\begin{equation}
P_{\rm obs}(\epsilon)=
\mathcal V P_{\rm ind}(\epsilon)
+(1-\mathcal V)P_{\rm dist}(\epsilon)+B ,
\label{eq:visibility-mixture}
\end{equation}
where $\mathcal V$ is a run-level mixture weight, $P_{\rm dist}$ is the
fully distinguishable labelled-photon probability in the same network, and
$B$ is an additive accepted-trial background.  Equation
\eqref{eq:visibility-mixture} is not asserted to be a universal
partial-distinguishability model; general Gram-matrix treatments are given
in Refs.~\cite{Shchesnovich2015,DufourBuchleitner2026}.

At $\epsilon=0.1$, $\mathcal V=0.99$, and $B=0$, the
$(0,1,2,1)$ target has
\begin{align}
P_{\rm obs}^{X_{13}}&=2.91093\times10^{-3},\nonumber\\
P_{\rm obs}^{Y_{13}}&=4.74917\times10^{-4}.
\end{align}
The exact ideal contributions are $2.46684\times10^{-3}$ and zero.
Under independent binomial sampling, the normal approximation for the
difference of two proportions gives
\begin{equation}
M\simeq
\frac{Z^2[P_X(1-P_X)+P_Y(1-P_Y)]}{(P_X-P_Y)^2},
\label{eq:shot-count}
\end{equation}
or $M=14228$ accepted four-photon trials per setting for $Z=5$.
This number is illustrative, model-dependent, and reproducible from the
released script; it is not a hardware-independent sample-complexity theorem.

The same model makes several failure modes explicit.  A small axis
misalignment $\delta$ away from $Y_{13}$ toward $X_{13}$ produces the
leading ideal term
\begin{equation}
P(\epsilon,\delta)=\frac{\epsilon^2\sin^2\delta}{4}
+o(\epsilon^2),
\end{equation}
so the null is a sensitive quadrature-alignment check.  Pure attenuation
after the interferometer multiplies a fixed target amplitude and cannot lift
an exact zero when conditioning on all $n$ detected photons.  By contrast,
internal or path-dependent loss can alter relative paths; source
contamination and detector noise can add a floor.  These contributions must
be independently characterized before counts on a null axis are assigned to
one cause.

\section{Relation to prior robustness results}

Dittel \emph{et al.} derived generic quadratic leakage under weak random
unitary disorder and linear leakage in small distinguishability parameters
after disorder averaging \cite{Dittel2018PRL}.  Our result is compatible
with that scaling but retains the coherent generator before averaging.
The response may vanish on a tangent direction, begin at fourth order, or
remain zero along an entire one-parameter subgroup.  Averaging over generic
directions suppresses precisely this anisotropic information.

Our four-bin histogram should also not be confused with the Fourier
transform over the permutation group used to resolve exchange-symmetry
sectors \cite{DufourBuchleitner2026}.  Here the bins count root-of-unity
phases of permanent paths in a fixed $F_4$ gauge.  The Krawtchouk step is a
sectorwise, phase-resolved application of mature symmetric-power machinery,
not a new Krawtchouk identity.

\section{Discussion}

The main conceptual lesson is that ``robustness of a suppressed event'' is
not a scalar without a perturbation model.  Selected $F_4$ events have
distinct response jets across experimentally addressable mixer
quadratures.  The cleanest near-term test uses four photons and the exact
$X_{13}/Y_{13}$ contrast in Eq.~\eqref{eq:headline}.  The
eleven-photon quartic direction is a higher-resource prediction.

Several limitations define the next work.  A hardware claim requires a
reconstructed circuit and a full distinguishability matrix rather than the
mixture in Eq.~\eqref{eq:visibility-mixture}.  Internal component errors
must be pushed through the circuit to determine which output tangent
directions they excite.  The finite closure census does not classify all
Fourier zeros, and the naive parity-sector analogue already fails for
$F_8$.  Finally, the sectorwise reciprocity lies close to classical
multivariate Krawtchouk theory; the priority claim made here is limited to
its $F_4$ phase-resolved use and the coherent response consequences.

\begin{acknowledgments}
The author used OpenAI Codex as a research and writing assistant for
symbolic exploration, code generation, literature discovery, and editorial
revision.  Every mathematical claim and numerical output in the manuscript
was checked against the cited sources or the released exact-arithmetic
programs; responsibility for the content remains with the author.
\end{acknowledgments}

\section*{Data and code availability}

All data are generated by the source code accompanying this manuscript.
The exact phase-histogram, reciprocity-census, leakage, and finite-shot
scripts use only the Python standard library.  The ancillary submission
archive contains the complete code and tests needed for the reported results.

\bibliography{references}

\end{document}
